\section{Methodology}
\label{sec:methodology}

GAME-Mal is a three-stage pipeline: (i) reflection-aware parsing of
Droidmon JSONL traces into API-call sequences, (ii) Markov-style
feature extraction used by our statistical baseline, and (iii) a
gated multi-head attention transformer that consumes the tokenised
sequences and emits both a family prediction and a per-token
importance distribution.

\subsection{Reflection-aware API resolution}
\label{sec:reflection}

Droidmon records every framework-level API call made by a running
APK as a JSON object with fields \texttt{class}, \texttt{method},
and — for calls that arrive through \texttt{java.lang.reflect} —
additional \texttt{hooked\_class} and \texttt{hooked\_method}
fields. A naive tokeniser emitting \texttt{class.method} would
collapse every reflective call in the trace to the generic token
\texttt{java.lang.reflect.Method.invoke}, destroying the
family-defining signal that malware authors hide behind reflection.

We therefore apply the following rule when converting a record to a
token $t$:
\[
  t =
  \begin{cases}
    \texttt{REFL:}\,h_c.h_m & \text{if } h_c,\, h_m \text{ present},\\
    c.m & \text{otherwise},
  \end{cases}
\]
where $c,m$ are the direct class and method and $h_c,h_m$ are the
hooked (true) class and method. The \texttt{REFL:} prefix is retained
deliberately: it is itself a behavioural feature — several families
(e.g.\ DroidKungFu) invoke reflection at characteristic rates.

A per-sample sequence $S = (t_1, t_2, \ldots, t_L)$ is the
time-ordered list of these tokens. We cap $L$ at $L_{\max}=512$ via
\emph{head}-truncation, i.e.\ keeping $S_{1{:}L_{\max}}$. We initially
hypothesised that tail-truncation would be preferable because it
preserves the payload-execution phase, but the seq-length sweep
(Sec.~\ref{sec:results}) decisively showed the opposite: head-truncation
beats tail-truncation by 1.7--2.9 macro-F1 points at every length
$\in\{256, 512, 768\}$. We attribute this to the early registration
and reflection-resolution events being more class-discriminative on
this corpus than the late command-and-control activity.

\subsection{Vocabulary and encoding}

Tokens occurring in fewer than $\kappa=2$ training samples are
mapped to a single \texttt{<UNK>} symbol, together with reserved
\texttt{<PAD>} and \texttt{<CLS>} tokens. The resulting vocabulary
contains $|V|=1{,}118$ symbols on our corpus. Each sequence is
prefixed with \texttt{<CLS>}, integer-encoded, and right-padded to
$L_{\max}$.

\subsection{Markov associative-rule features}
\label{sec:markov}

Following D'Angelo et al.~\cite{dangelo2023association} we treat each
sequence as a source for \emph{$k$-spaced associative rules}. For
every ordered pair $(A,B)$ of tokens with a gap $k \in \{1,\ldots,K\}$
(we use $K=10$), we emit the rule $A \xrightarrow{k} B$ and count
its occurrences within and across samples.

For each rule $R$ we compute:
\begin{align}
  \mathrm{supp}(R) &= \frac{\#\{\text{samples containing }R\}}{N},\\
  \mathrm{conf}(R \mid c) &= \frac{\#\{\text{samples of class }c\text{ containing }R\}}{\#\{\text{samples containing }R\}},
\end{align}
where $N$ is the number of training samples. We keep a rule if
$\mathrm{supp}(R) \geq 5\times 10^{-4}$ and
$\max_c \mathrm{conf}(R \mid c) \geq 0.30$. On our corpus this
prunes $\sim\!90\%$ of the rule space, leaving roughly 2{,}200 rules
per fold. For the rule-based classifier we implement D'Angelo's
Eq.~6 faithfully:
\begin{equation}
  \rho_c(S) = \sum_{R \in \mathcal{R}(S)}
  \frac{\sigma_S(R)}{|S|}\,\gamma(R, c),
  \label{eq:rho}
\end{equation}
where $\sigma_S(R)$ is the number of occurrences of rule $R$ in
sequence $S$ and $\gamma(R, c) = \mathrm{conf}(R\mid c)$. We
normalise $\rho$ via softmax to obtain class probabilities and
predict $\hat{c} = \arg\max_c \rho_c$.

\subsection{GAME-Mal architecture}
\label{sec:arch}

The learned component is a pre-norm transformer encoder with
$N=2$ layers, $d_{\text{model}}=128$, $h=4$ heads, and a
feed-forward width of $d_{\text{ff}}=256$. Dropout is set to $0.1$
throughout. The total parameter count is 515{,}656.

\paragraph{Embedding and positional encoding.} Tokens are embedded
via a learned lookup $E \in \mathbb{R}^{|V|\times d_{\text{model}}}$,
followed by sinusoidal positional
encodings~\cite{vaswani2017attention}. Let $X \in
\mathbb{R}^{L_{\max}\times d_{\text{model}}}$ be the resulting
sequence.

\paragraph{Gated multi-head attention (G1).} Following
Qiu et al.~\cite{qiu2025gated}, we compute standard scaled
dot-product attention per head $i\in\{1,\ldots,h\}$:
\begin{equation}
  A_i = \mathrm{softmax}\!\left(\frac{Q_i K_i^\top}{\sqrt{d_k}}\right) V_i,
  \quad d_k = d_{\text{model}}/h,
\end{equation}
and then \emph{gate} each head's output with an input-dependent
sigmoid before the output projection $W_O$:
\begin{align}
  g_i &= \sigma(X W_{g,i} + b_{g,i}),\quad g_i\in\mathbb{R}^{L_{\max}\times d_k},\\
  A_i' &= g_i \odot A_i,\\
  Y &= \mathrm{concat}(A_1',\ldots,A_h')\, W_O.
  \label{eq:g1}
\end{align}
The gate bias $b_{g,i}$ is initialised to $-2.0$
($\sigma(-2)\approx 0.12$) following Qiu et al.; the final learned
distribution on our task is considerably less sparse
(Section~\ref{sec:results}). The gate acts as a per-token,
per-head input-dependent scaling: tokens whose $g_i$ is near zero
are suppressed before the output projection, which both breaks the
$V$--$W_O$ low-rank bottleneck and eliminates the attention-sink
artefact empirically.

\paragraph{Feed-forward block.} A standard two-layer GELU MLP with
residual connections and pre-layer normalisation follows each
attention block.

\paragraph{Classification head.} We apply mean pooling over
non-pad positions, then a single linear layer to the $|C|=8$
family logits. Classification uses a class-weighted cross-entropy
loss
\begin{equation}
  \mathcal{L} = -\sum_{i=1}^{B} w_{y_i}\log \mathrm{softmax}(z_i)_{y_i},
\end{equation}
with $w_c = N / (|C|\,N_c)$, partially compensating for the
$34\times$ imbalance between the largest and smallest families.

\subsection{Training protocol}
\label{sec:training}

Models are trained with AdamW
($\beta_1{=}0.9$, $\beta_2{=}0.999$, weight decay $10^{-4}$), a peak
learning rate of $3\times 10^{-4}$ under a cosine annealing schedule
with 5-epoch linear warm-up, gradient clipping at $\|g\|_2=1.0$, and
a batch size of 32. We train for up to 40 epochs with early stopping
on the validation macro-F1 (patience of 10 epochs). Runs use mixed
\texttt{float32} on the Apple Metal Performance Shaders (MPS)
backend of PyTorch; all randomness is seeded per fold.

\subsection{Intrinsic explainability}
\label{sec:explain-method}

Because Eq.~\ref{eq:g1} lies on the forward path of the classifier,
the gate activations $g_i^{(\ell)}(x)$ at layer $\ell$, head $i$, for
input token $x$ are themselves an explanation: they quantify how
much of that token's contribution survives the bottleneck on its way
to the classifier. We aggregate
$\bar{g}(x) = \frac{1}{hN}\sum_{\ell,i} g_i^{(\ell)}(x)$
over heads and layers and, for every family, report the top-$k$
APIs ranked by mean $\bar{g}$ over true-positive samples of that
family. This produces the per-family behavioural fingerprints of
Section~\ref{sec:explainability} without any post-hoc surrogate.
